\chapter{Definitions}

The following terms are used in this work:\par\bigskip

\begin{description}[style=multiline, labelwidth=4.2cm, leftmargin=4.6cm]

  \item[external penetration testing]
       A security assessment performed from the public side of the network, without access to source code, credentials, or internal infrastructure of the target. The tester interacts only with publicly reachable endpoints.

  \item[black-box testing]
       A testing approach in which the assessor has no prior knowledge of the internal structure of the target system; conclusions are based exclusively on observed external behavior.

  \item[reconnaissance]
       The phase of a penetration test devoted to collecting public information about a target: domain and DNS records, subdomains, hosted services, technologies, exposed files, and historical artifacts.

  \item[active probe]
       A scanning component that issues purposeful HTTP requests to a target with the intent of triggering a measurable side effect (reflection, error, timing deviation, redirect) that reveals a security weakness.

  \item[passive analysis]
       A scanning component that derives security findings without sending additional requests to the target, typically by inspecting responses, headers, JavaScript files, or third-party data sources.

  \item[finding]
       A normalized record describing a single security issue or candidate issue, with a stable identifier, severity, confidence, evidence, and references to CWE, OWASP, and MITRE ATT\&CK.

  \item[verdict]
       A triage label attached to a finding (\textit{confirmed}, \textit{candidate}, or \textit{manual\_review}) that indicates how much human verification is still required before the finding can be acted upon.

  \item[scope]
       The set of domains, subdomains, and IP addresses that the operator has authorized the framework to test. All HTTP traffic that would target hosts outside this set is blocked by the framework before being sent.

  \item[scan profile]
       A predefined set of configuration overrides (\textit{safe}, \textit{bug\_bounty}, \textit{deep}, \textit{intrusive}) that enables or disables intrusive behaviors such as profile comparison in IDOR testing or active confirmation in race-condition probes.

  \item[Large Language Model (LLM)]
       A neural network trained on a large text corpus that is capable of producing coherent natural-language explanations from a structured prompt. In this work, an LLM is used post-scan to generate an analyst-style report from the raw findings.

\end{description}
